\section{Introduction to the Topic and Historical Context}

\lead{Throughout Latin America in 1975, the political future of the continent remains uncertain. Military governments, civilian administrations, and political and military authorities confront growing instability, ideological porlarization, and the possibility that developments beyond their borders may affect their own political survival.}

The expansion of revolutionary movements following the Cuban Revolution has contributed to a growing perception that political conflict can no longer be contained within national borders. Revolutionary organizations operate across countries, political opponents seek refuge abroad, and exile communities provide new spaces for political organization. In response, military and intelligence institutions have gained increasing influence over domestic politics.

These events have made governments across the region face an increasingly interconnected political environment. Developments in neighboring countries, the interests of external powers, and domestic instability may all affect their position.

\subsection{Key Terms}

\begin{guidetable}[Y{0.45}Y{1.55}]{2}
  \textbf{Southern Cone} & The most southern countries within South America (Argentina, Chile, Uruguay, Paraguay and sometimes Brazil/Bolivia). The term denotes both a geographical area and a shared political trajectory: by the 1970s, most of these states were under military rule pursuing near-identical anti-communist security agendas (Anderson, 2023). \\
  \midrule
  \textbf{Junta} & Usually a military group that seizes and holds governing power outside constitutional processes, oftentimes following a coup. In the context regarding the Southern Cone, juntas present themselves as temporary custodians of national order instead of conventional governments, which shape how they justify extraordinary security measures (Merriam-Webster, 2019). \\
  \midrule
  \textbf{National Security Doctrine} & From the 1950s onward, this framework has been developed across Latin American militaries aiming at redefining internal political dissent as a security threat equivalent to military invasion. Under this doctrine, unions, student groups, political parties and the press can all be classified as fronts for subversion, which give militaries a rationale to intervene in civilian life far beyond conventional defense functions (Burning Issues, 2025). \\
  \midrule
  \textbf{Bureaucratic-} \newline \textbf{authoritarianism} & A specific style of military rule that emerged in the Southern Cone: technocratic, rules-based and organised around economic stabilisation and anti-communist order, rather than personality-driven dictatorship (Collier, 2001) \\
  \midrule
  \textbf{DINA} & Chile's secret intelligence police under Pinochet, created in 1973 and directed by Manuel Contreras. DINA operates with near-total autonomy from civilian oversight and has become the main architect of Chile's international repression network, including operations well beyond Chile's borders (Explaining History, 2025). \\
  \midrule
  \textbf{SIDE} & Argentina's state intelligence secretariat, which plays a parallel role to DINA within Argentina's own security apparatus (Grok, 2026). \\
  \midrule
  \textbf{School of the Americas} & A US Army training institution (based in Panama) in which thousands of Latin American military officers receive instruction in counterinsurgency, intelligence and related doctrine. It functions as one of the key channels through which a shared anti-communist doctrine and shared professional networks spread across the region's officer corps (Grimmett \& Sullivan, 2001). \\
  \midrule
  \textbf{Domino Theory} & The Cold War-era belief, held primarily by US policymakers, that if one state in a region fell to communism, its neighbours would follow in a chain reaction. This logic is used to justify US intervention and support for anti-communist regimes throughout Latin America, regardless of how moderate the actual political threat is (Britannica, 2018). \\
  \midrule
  \textbf{Alliance for Progress} & A US economic and technical assistance program launched in 1961, intended to counter the appeal of the Cuban Revolution by promoting development and reform in Latin America. In practice, it has become closely tied to US security objectives in the region (John F. Kennedy Presidential Library and Museum, 2019). \\
\end{guidetable}

\subsection{The Cold war in the Southern Cone}

\subsubsection{National Security Doctrine}

The Southern Cone was the frontline of the Cold War by the late 1960s. US national security planners and regional elites had fused anti-communism with older, homegrown fears of social revolution (McSherry, 1999, 2007). Security establishment in Argentina, Chile, Uruguay, Paraguay and Brazil started treating domestic opponents, trade unionists, students, Christian democrats, left parties, even moderate reformers and 'internal enemies’ aligned with an international subversive conspiracy, rather than as legitimate players in pluralist politics (McSherry, 2007; Zanchetta, 2018). That reframing sets the stage for the counter-revolutionary system which authoritarian regimes adopted. A system that presupposed elites who have already decided their political opponents are foreign agents in disguise (Lessa \& Balardini, 2024; McSherry, 2005).

From Washington’s side, Latin America was seemingly a vulnerable player. Once Cuba aligned itself with the Soviet bloc after 1959, US strategists considered even modest redistributive or nationalism projects in the Southern Cone as potential dominoes, places where a ‘loss’ could tip a much wider geopolitical balance (McSherry, 2007; Zanchetta, 2018). Washington paired this anxiety with inducement, introducing the Alliance for Progress. Launched in 1961, the Alliance offered economic and technical assistance meant to counter the appeal of the Cuban Revolution by promoting development and land reform through legitimate channels rather than armed struggle. In practice, the program’s implementation became increasingly bound up with the same security logic driving containment policy, so that development aid and counterinsurgency support often flowed through the same institutions and towards the same recipients (Taffet, 2007). Southern Cone militaries and conservative civilians took that logic and executed it by introducing agrarian reforms, stronger labour rights and resource nationalism, which were all interpreted as symptoms of incipient insurgency, not as normal democratic contestation (McSherry, 1999, 2007).

The clearest early example of this dynamic emerged in Brazil, where the 1964 military coup built a security apparatus explicitly organised around the ‘Doctrine of National Security and Development’, a framework taught at the Escola Superior de Guerra that redefined the military’s role as guardian against international subversion rather than external invasion. Brazilian officers exported this model through joint training exercises with Argentine, Chilean, Uruguayan and Paraguayan counterparts as early as 1969-70, effectively acting as an early laboratory for cross-border professional relationships(McSherry, 1999).

This way, ‘national security states’ were consolidating across the region with militaries claiming a near-messianic mandate to remake society through anti-communism, building out clandestine intelligence structures and parallel forces under direct US guidance (McSherry, 2005, 2007; Huggins, 1998). Much of this professionalisation ran through shared inter-American training pipelines, most notably the US Army’s School of the Americas and the network of national war academies it fed into, where officers from different countries absorbed the same doctrine, the same enemy list and increasingly the same operational habits (McSherry, 2005).

Central to this doctrine was the idea that the ‘enemy’ was defined by ideas and loyalties, not by a specific criminal act. This premise, refined in shared training centres and circulated as doctrine across the region, is what let security concerns go over into every aspect of civilian life, unions, universities, churches, the press, largely without institutional resistance (McSherry, 1999, 2007; Lessa \& Balardini, 2024). It also explains why much of the repression  had a distinctly bureaucratic, procedural character: officers were not improvising against a visible threat but executing a doctrine that had already defined the shape of the enemy in advance. By the early 1970s, authoritarian security services already spoke the same doctrinal language, used the same operational repertoire and pointed at the same enemy image, long before they organised formally around a shared name (McSherry, 2005; Lessa \& Balardini 2024).

\subsubsection{The Cuban Effect \& Rise of Guerrilla Movements}

The 1959 Cuban Revolution did not only install a socialist government about ninety miles from Florida, it also supplied an entire generation of the Latin American left with a working model of what armed struggle could achieve. Throughout the 1960s, movements explicitly inspired by Havana’s example emerged across the continent such as Argentina’s Monotoneros and ERP, Uruguay’s Tupamaros, Chile’s MIR, alongside sister organisations in Colombia, Venezuela and Bolivia (Brown, 2017). Most of these groups had no direct organisational link to Cuba and had never trained there, but they shared its basic premise, that electoral politics was a dead end and that armed vanguard action could substitute for the slow work of mass mobilisation (Brown, 2017).

This was not merely rhetorical inspiration as Cuban intelligence, operation through what later became known as the Departamento América, maintained real if uneven contact with various Latin American armed groups, offering training, safe passage and occasional material support, under a stated policy of persuasion rather than direct command and control (Domínguez, 1989). Che Guevara’s own attempt to export the Cuban model to Bolivia became a cautionary tale on both sides, proving to guerillas that revolution could travel while showing to regional militaries that the threat was not hypothetical.

For the aforementioned security establishments, this wave of urban guerilla activity was less a discrete security problem than confirmation of an existing worldview. Where the National Security Doctrine had already primed officers to see internal dissent as evidence of a broader subversive conspiracy, the actual appearance of armed, ideologically explicit, Cuban-inspired movements closed the loop. The guerillas were treated not as a domestic law and order issue but as the local expression of a hemispheric war already underway. That framing is precisely what makes a transnational, rather than a purely national response appear both necessary and self-evidently justified (McSherry, 1999, 2007).

\subsubsection{National Case Studies}

\minorheading{Argentina}

Argentina entered the 1970s already deeply into a cycle of instability that predated any coordinated regional response. Juan Perón’s return from exile in 1973, after nearly two decades of proscription, briefly appeared to offer a path back to civilian consensus. Instead, the Ezeiza massacre that greeted his return exposed the irreconcilable split between the Peronist left and right and Perón’s own subsequent break with the Montoneros accelerated rather than resolved the conflict (Gillespie, 1982). Following Perón's death in 1974, his widow Isabel Perón inherited a state that was losing control on two fronts simultaneously with left-wing guerilla organisations, specifically the Montoneros and the Marxist ERP, escalating kidnappings and assassinations of military and police officers, while the right-wing Argentina Anticommunist Alliance, known as the Triple A, organised with the tacit backing of figures inside her own government, conducted its own campaign of killings against trade unionist and left-wing Peronists (Lewis, 2002). By 1975, this had become state policy in substance with Isabel Perón's government issuing decrees explicitly empowering the military to annihilate left-wing subversive elements as part of Operativo Independencia, amid inflation running above 300 percent that compounded a sense of a government losing its grip on both the economy and the security situation (Lewis, 2002). Argentina thus arrived at the threshold of formal regional coordination already primed by its own domestic emergency, with a security apparatus increasingly willing to operate outside conventional legal constraints.

\minorheading{Chile}

Chile’s trajectory was more compressed but equally formative with Salvador Allende’s 1970 election as the first democratically elected Marxist head of state in the hemisphere, triggering immediate and sustained US concealed opposition from the Nixon administration, instructing to make the economy suffer through years of financial destabilisation and supporting opposition media and unions (Harmer, 2011; Kornbluh, 2003). Domestic polarisation, driven both by external pressure and by internal contradictions within Allende’s own Popular Unity coalition, culminated in the September 1973 coup that installed Augusto Pinochet. What followed was unusually fast institutionalisation of repression as Pinochet secretly established what became DINA within weeks of the coup, formalising it by decree in June 1974. DINA operated entirely outside the ordinary military chain of command, reporting directly to Pinochet alone (Kornbluh, 2003). Within a year, DINA had already demonstrated a willingness to act beyond Chile’s borders, carrying out the September 1974 car bomb assassination of exiled former army commander Carlos Prats in Buenos Aires, an operation frequently seen as the first instance of the Chilean state exporting political violence abroad (Kornbluh, 2003). Chile's case therefore shows the most direct line from domestic coup to cross-border operational capacity in the Southern Cone.

\minorheading{Uruguay}

Uruguay’s slide followed a slower but equally consistent logic as it has long been considered the most stable democracy in the region. The country’s institutions began eroding under President Jorge Pacheco Areco, who from 1968 onward repeatedly invoked emergency prompt security measures to suspend civil liberties in response to the Tupamaros, a Marxist urban guerilla group known for high profile kidnapping and bank robberies (Weinstein, 1988). By 1972, under Pacheco’s successor Juan María Bordaberry, the state had declared a formal state of internal war, granting the military expanded powers of arrest and interrogation. Through sustained use of torture and mass detention, the armed forces had substantially dismantled the Tupamaros’ leadership by the end of that year (Weinstein, 1988; Roniger \& Sznaider, 1999). Notably, this meant that guerrilla threat that had originally justified military involvement in Uruguayan security policy was already largely defeated well before the June 1973 coup, in which Bordaberry, with military backing, dissolved parliament altogether and formalised military control through a new National Security Council (Roniger \& Sznaider, 1999; Lessa, 2022). Uruguay’s case is therefore something of an outlier within the group, with military and security consolidation continuing to deepen even after its stated jurisdiction had largely disappeared, suggesting that by the mid-1970s the security apparatus itself, not merely the threat it was built to counter, has become the organising logic of the state.

\subsubsection{Precursors to Condor (1969 - 1975)}

Before any formal agreement existed, Southern Cone security services had already begun sharing intelligence on an ad hoc basis. Chilean, Paraguayan, Argentine and Uruguayan agencies exchanged information on subversives informally throughout the early 1970s, largely through personal contacts between officers rather than through any standing institutional channel (Dinges, 2004). This informality mattered as not cooperation depended on individual relationships and mutual trust built up over years, rather than on any binding commitment.

As early as 1969 and 1970, Brazilian military bases hosted counterinsurgency training for Argentine, Chilean, Uruguayan and Paraguayan officers, bringing together the personnel who operated coordinated operations (Huggins,1998). Joint teams formed during this period gathered intelligence that was subsequently used in domestic political repression, establishing a working precedent for cross-border collaboration (Huggins, 1998; McSherry, 2005).

The 1973 coups in Chile and Uruguay marked the point at which this informal cooperation became urgent rather than merely convenient. With both countries now under military rule and actively hunting the same category of political opponents, often the same individuals, as they crossed borders into exile, the existing informal networks have proved too slow and too unreliable for the scale of coordination the juntas want (Lessa, 2022). What had been a loose set of bilateral courtesies now needed to become something closer to a system.
